\section{个人成长与教育情况}
\subsection{团队学习成长路径}

项目以完整推理任务牵引学习，将指令执行、存储、操作系统、量化计算和软件接口联系起来。成员从解释模块原理，逐步转向分析接口约束、定位联合问题和设计实验。学习成效通过设计复述、代码审查、失败用例分析及独立复现记录，体现专业知识解决实际问题的过程。\cite{reviewrules}

\subsection{研发历程}

当前已形成\TechnicalBaseline{}技术选型与可行性报告，完成资料、源码接口和主要依赖许可核查；RTL生成、软件构建、仿真、综合、系统移植、模型训练及上板验证尚未开展。后续以此为起点，记录问题、设计判断、测试结果和修订原因，将资料核查与工程验证分别说明。\cite{technicalreport}

\subsection{教育历程}

团队计划补足三方面能力：体系结构与FPGA知识用于理解处理器、向量单元及矩阵加速器接口；操作系统知识用于启动、内存、线程上下文和驱动适配；机器学习与数值计算知识用于模型准备、量化语义和误差分析。课程、实验及科研训练经历按成员实际情况补充，以能够解释和验证的工作体现能力。

\subsection{产教融合}

围绕\ChallengeCompany{}的命题，将要求拆分为功能、接口和验证条件。企业对接重点为国产芯片实现口径、演示场景及移动数据中心接口；教师指导方案可行性与研究方法；学生负责设计、实现、验证及记录。实际反馈进入需求响应表并推动方案修订，不预设企业已经认可方案或提供合作证明。

\subsection{未来规划}

首先准备模型资产与独立整数参考，验证受限算子及三后端计算；随后推进SoC配置、RTL仿真、资源和时序检查，再开展板上串口、DDR与启动适配。两种系统分别启动后，核对向量上下文、NPU搬运和重复推理。各阶段以可复核交付物验收，资源或时序未通过时先收敛配置。

基础推理通过后，再依据企业反馈选择静态视觉展示和感知通信扩展。责任人、时间与验收记录随分工更新，持续区分设计目标和实际结果，使成长记录对应真实问题与解决过程。
